\cleardoublepage
\chapternonum{作者简历}

\section*{基本信息}

翁玮，女，1999年4月生，浙江省人。
2017年9月进入浙江大学环境与资源学院资源环境科学专业学习，攻读学士学位。
2021年9月继续在浙江大学环境与资源学院农业遥感与信息技术专业直接攻读博士学位，师从黄敬峰教授。
博士期间主要从事农业遥感与气候变化研究，研究内容包括全球水稻分布制图、全球水稻种植格局变化及其气候效应等。

\section*{攻读博士学位期间完成的论文}

\begin{enumerate}
    \item \textbf{Weng, W.}, Huang, J., Yue, C., Lyu, Z., Xiao, Y., Li, S., Huang, R., Peng, D., Huang, C., Bigirimana, F., Liu, L., Liu, W., 2026. Widespread land surface cooling from paddy rice cultivation revealed by global satellite mapping. \textit{\textbf{Nature Communications}}. 17, 844. https://doi.org/10.1038/s41467-025-67549-z. (SCI, 5-year IF = 18.9, 2025 JCR)
    
    \item Xiao, Y., Huang, J., \textbf{Weng, W.}, Huang, R., Shao, Q., Zhou, C., Li, S., 2024. Class imbalance: A crucial factor affecting the performance of tea plantations mapping by machine learning. \textit{\textbf{International Journal of Applied Earth Observation and Geoinformation}}. 129, 103849. https://doi.org/10.1016/j.jag.2024.103849. (SCI, 5-year IF = 9.1, 2025 JCR)

    \item Huang, R., Xiao, Y., Li, S., Li, J., \textbf{Weng, W.}, Shao, Q., Zhang, J., Zhang, Y., Yang, L., Huang, C., Sun, W., Liu, W., Jin, H., Huang, J., 2025. A novel framework for dynamic and quantitative mapping of damage severity due to compound drought--heatwave impacts on tea plantations, integrating Sentinel-2 and UAV images. \textit{\textbf{Computers and Electronics in Agriculture}}. 228, 109688. https://doi.org/10.1016/j.compag.2024.109688. (SCI, 5-year IF = 10.7, 2025 JCR)
    
    \item Xiao, Y., Zhao, Z., Huang, J., Huang, R., \textbf{Weng, W.}, Liang, G., Zhou, C., Shao, Q., Tian, Q., 2024. The illusion of success: Test set disproportion causes inflated accuracy in remote sensing mapping research. \textit{\textbf{International Journal of Applied Earth Observation and Geoinformation}}. 135, 104256. https://doi.org/10.1016/j.jag.2024.104256. (SCI, 5-year IF = 9.1, 2025 JCR)
    
    \item Huang, R., Li, S., Zhu, X., Li, J., Xiao, Y., \textbf{Weng, W.}, Shao, Q., Chai, D., Zhang, J., Zhang, Y., Yang, L., Wu, K., Hu, Z., Liu, L., Sun, W., Liu, W., Huang, J., 2025. A novel scheme for seamless global mapping of daily mean air temperature (SGM\_DMAT) at 1-km spatial resolution using satellite and auxiliary data. \textit{\textbf{Ecological Informatics}}. 90, 103266. https://doi.org/10.1016/j.ecoinf.2025.103266. (SCI, 5-year IF = 8.0, 2025 JCR)    
    
    \item Zhao, Z., Wang, F., Xiao, Y., Liu, L., Chen, N., \textbf{Weng, W.}, Peng, D., Zhang, X., Shi, Z., 2025. Development and evaluation of remote sensing-derived relative growth rate (RGR) for rice yield estimation. \textit{\textbf{IEEE Journal of Selected Topics in Applied Earth Observations and Remote Sensing}}. 18, 25735--25747. https://doi.org/10.1109/JSTARS.2025.3617020. (SCI, 5-year IF = 6.2, 2025 JCR)
    
    \item Shao, Q., Huang, C., Xiao, Y., Liu, L., Liu, W., Huang, R., Zhou, C., \textbf{Weng, W.}, Huang, J., 2025. A background field-based framework for reconstructing spatially and temporally continuous daily NDVI (1982--2020) from the AVHRR LTDR. \textit{\textbf{Ecological Indicators}}. 181, 114441. https://doi.org/10.1016/j.ecolind.2025.114441. (SCI, 5-year IF = 8.5, 2025 JCR)
    
    \item Shao, Q., Huang, C., Xiao, Y., Liu, L., Liu, W., Huang, R., Zhou, C., \textbf{Weng, W.}, Huang, J., 2025. Selecting of global phenological field observations for validating coarse AVHRR-derived forest phenology products based on spatial heterogeneity and temporal consistency. \textit{\textbf{Ecological Informatics}}. 90, 103216. https://doi.org/10.1016/j.ecoinf.2025.103216. (SCI, 5-year IF = 8.0, 2025 JCR)
    
    \item Xiao, Y., Huang, R., Liu, L., Zhu, L., Shi, Z., Zhou, C., \textbf{Weng, W.}, Sun, W., Liu, W., Shao, Q., Zhao, Z., Huang, J., 2026. A harmonized 10 m near-2-day optical time series dataset improves tea plantation mapping and pruning date estimation. \textit{投稿中}.
\end{enumerate}


\section*{攻读博士学位期间参与的科研项目}

\begin{enumerate}
    \item 耦合先验知识与深度学习的水稻种植区遥感提取方法研究 [面上项目]
    \item 作物涝渍灾害立体动态监测技术与智能评估系统研发 [国家重点研发子课题]
    \item 长江中下游地区水稻主要气象灾害监测技术方法研究 [国家重点研发子课题]
    \item 地理信息系统和遥感技术在林业与生态领域的可持续应用 [欧盟国际合作]
    \item 茶叶气象监测关键技术研究与靶向服务 [浙江省重点研发子课题]
    \item 河北省气象灾害防御信息化平台建设 [横向]
\end{enumerate}
